\chapter{TRABAJOS PREVIOS} \label{chap:trabajos-previos}

En los últimos años, las técnicas de aprendizaje profundo han demostrado un rendimiento
excepcional en numerosas tareas que van desde el reconocimiento de objetos hasta el
procesamiento del lenguaje natural. El uso de técnicas de aprendizaje profundo,
específicamente redes neuronales convolucionales (CNN), ha demostrado resultados
prometedores en la detección automática de RD a través del análisis de imágenes del fondo
de ojo. Dentro de este contexto, se incluyen varios enfoques que han contribuido en este
campo de investigación.

\section{Enfoques con CNN personalizadas}

Este problema se ha abordado de distintas formas, planteándose distintas soluciones. En
general, podemos ver en la literatura una preponderancia de la implementación de modelos
previamente entrenados. Esto toma sentido cuando consideramos el costo de entrenamiento
que tienen este tipo de modelos que se alimentan de cientos de miles de imágenes; solo
pocos sectores de la industria tienen los recursos necesarios para llevar a cabo estos
desarrollos. Muchos de los artículos recopilatorios encargados de comparar soluciones
frente a la clasificación de la RD se encargan, en última instancia, de comparar la
efectividad de distintos modelos preentrenados para la clasificación de la RD
\myfootcite{bala2024comparative}, donde se comparan ocho de los modelos preentrenados más
utilizados desde 2011 hasta 2022.

En algunas investigaciones se observa un enfoque común: la personalización de un modelo
de red neuronal convolucional (CNN) para abordar la detección automatizada de retinopatía
diabética (RD). Esta estrategia implica modificar la arquitectura o los parámetros de un
modelo CNN tradicional y mejorarlo en la tarea específica de la detección de RD. En este
contexto, se han explorado diversas estrategias para mejorar la precisión y eficiencia de
la detección de RD a partir de imágenes del fondo de la retina. Por ejemplo, se han
aplicado etapas de preprocesamiento como escalado, recorte y ecualización adaptativa del
histograma con limitación de contraste (CLAHE) para mejorar la calidad de las imágenes
\myfootcite{tabtaba2024diabetic}.

Además, se han utilizado técnicas de aumento de datos mediante redes generativas
adversariales (GAN) para enriquecer el conjunto de datos de entrenamiento. Otros trabajos
han introducido variantes personalizadas de CNN, como la Red Neuronal Convolucional
Dilatada Multiescala Híbrida en Cascada (HCMD-CNN) \myfootcite{gayathri2020lightweight},
que integra la Red de Atención Residual (RAN) y MobileNet, optimizando sus parámetros
para lograr un mejor rendimiento en la detección de RD. Otros enfoques han empleado
métodos de extracción de características mediante CNN, cuyas salidas se utilizan como
entrada para diferentes clasificadores de aprendizaje automático
\myfootcite{li2019computer} \myfootcite{gayathri2021diabetic}, como Support Vector
Machine, AdaBoost, Naive Bayes, Random Forest y J48.

\section{Enfoques con arquitecturas preentrenadas}

Últimamente se han desarrollado otras aproximaciones que aprovechan arquitecturas
preentrenadas en el ámbito de la detección y clasificación de retinopatía diabética (RD)
en imágenes del fondo de ojo. Estos enfoques abarcan una variedad de modelos de
aprendizaje profundo, tales como VGG16, InceptionResNetv2, Inceptionv3, DenseNet121,
MobileNetV2, Xception y YOLOv3, los cuales se han utilizado para extraer características
y llevar a cabo tareas de clasificación y detección.

Por ejemplo, un estudio implementó una red neuronal convolucional (CNN) utilizando VGG16
para la extracción de características de imágenes de retina y la técnica de sobremuestreo
de minoría sintética (SMOTE) para manejar el desequilibrio de clases en el conjunto de
datos \myfootcite{mustafa2024rapid}, logrando una precisión del 93.94\% durante la fase
de entrenamiento y del 88.19\% durante la validación.

Otro enfoque consistió en dos modelos basados en aprendizaje profundo: CNN512, que
clasificó imágenes completas en estadios de RD con una precisión del 88.6\%, y un modelo
YOLOv3 adaptado para detectar y localizar lesiones de DR, alcanzando un mAP de 0.216 en
la localización de lesiones \myfootcite{alyoubi2021diabetic}. Estos modelos se fusionaron
para clasificar imágenes de RD y localizar lesiones, obteniendo una precisión del 89\%
con sensibilidad del 89\% y especificidad del 97.3\%.

Además, se ha explorado el aprendizaje por transferencia utilizando modelos preexistentes
como InceptionResNetv2 e Inceptionv3, junto con un nuevo modelo DiaCNN, para la
clasificación de enfermedades oculares en el conjunto de datos ODIR
\myfootcite{shoaib2024deep}. Estos modelos alcanzaron impresionantes niveles de
precisión, llegando hasta el 99.7\% en algunas pruebas.

Asimismo, se ha abordado la detección automática de vasos sanguíneos retinianos
utilizando modelos como ResNet34, InceptionV3 y VGG16, y se han propuesto enfoques
multinivel \myfootcite{beghriche2023multi} para mejorar la clasificación binaria y
multiclase de imágenes de RD \myfootcite{vij2024hybrid}, logrando una precisión del
96.50\% en la clasificación binaria y una kappa ponderada cuadrática del 93.00\% en la
clasificación por etapas. Estos avances muestran mejoras significativas en precisión y
generalización, con resultados que superan a los métodos tradicionales y a los trabajos
previos en la literatura.

\section{Enfoques con modificaciones arquitectónicas innovadoras}

Por otra parte, existen otros trabajos en la literatura que han optado por realizar
modificaciones sustanciales a las arquitecturas CNN para obtener un mejor desempeño. Se
han propuesto varias técnicas innovadoras para mejorar el rendimiento y la eficiencia del
aprendizaje profundo en la detección y clasificación de retinopatía diabética (RD)
\myfootcite{patil2023study}.

Un estudio aborda el aprendizaje profundo distribuido (DDL) para entrenar modelos de DL
de manera rápida, utilizando dos GPU y analizando el rendimiento del canal de paralelismo
en DDL, logrando una puntuación F1 de validación de 0.834 en la clasificación de RD. Otro
enfoque se centra en la diversidad de los conjuntos de datos de entrenamiento mediante la
manipulación de datos de entrada, mejorando así el rendimiento de la clasificación
\myfootcite{patil2022pipeline}.

Además, se propone un enfoque bioinspirado en la metaplasticidad sináptica, que se
incorpora en redes neuronales convolucionales para fortalecer las ocurrencias menos
comunes durante el aprendizaje \myfootcite{vives2021diabetic}. Este método mejora tanto
la tasa de aprendizaje como la precisión, superando a otros métodos de la literatura
actual incluso con conjuntos de datos más pequeños, con resultados destacados en la
arquitectura InceptionV3.

\section{Limitaciones y brechas de investigación}

A pesar del avance significativo en el uso de técnicas de aprendizaje profundo para la
detección de RD a partir de imágenes del fondo de ojo, aún existen algunas limitaciones
que deben abordarse para garantizar su aplicabilidad en el mundo real y su adopción en la
práctica clínica.

La mayoría de los modelos actuales se entrenan con conjuntos de datos limitados y
específicos, lo que limita su capacidad para generalizar a escenarios del mundo real con
mayor variabilidad en las imágenes del fondo de ojo. En muchos conjuntos de datos, las
imágenes con RD representan una pequeña proporción en comparación con las imágenes
normales. Este desequilibrio de clases puede afectar el rendimiento del modelo, sesgando
la clasificación hacia la clase mayoritaria.

Algunos estudios se enfocan únicamente en la detección de la presencia o ausencia de RD,
sin considerar la clasificación de su gravedad en estadios tempranos, moderados o
severos, lo cual es crucial para la toma de decisiones oportunas en el manejo de la
enfermedad. El ruido, la iluminación inconsistente y el bajo contraste en las imágenes
del fondo de ojo pueden afectar significativamente la precisión de la clasificación. La
implementación de técnicas de preprocesamiento de imágenes robustas es fundamental para
mejorar el rendimiento de los modelos.

Pocos trabajos han explorado sistemáticamente el efecto del procesamiento de canales de
color individuales (rojo, verde, azul) en el desempeño de clasificación de RD, a pesar de
que diferentes canales pueden resaltar lesiones específicas. Además, existe una escasez
de validación externa rigurosa en múltiples datasets con características de adquisición
diferentes, lo cual es esencial para demostrar la capacidad de generalización de los
modelos.

\section{Contribución del presente trabajo}

El empleo de algoritmos de aprendizaje profundo en la detección de RD a partir de
imágenes del fondo de ojo presenta un gran potencial para mejorar la accesibilidad, la
precisión y la eficiencia del tamizaje de esta enfermedad. Sin embargo, es necesario
abordar las limitaciones actuales relacionadas con la generalización de los modelos, el
desequilibrio de clases, la clasificación de la gravedad y el preprocesamiento de
imágenes para garantizar su aplicabilidad en entornos clínicos reales y su adopción como
una herramienta confiable para el diagnóstico y seguimiento de la RD.

El presente trabajo contribuye al estado del arte en los siguientes aspectos:

\begin{itemize}
    \item \textbf{Análisis sistemático de canales de color:} Evaluación comparativa del
desempeño de modelos procesando imágenes RGB completas versus canales individuales (rojo,
verde, azul), identificando el canal más efectivo para clasificación multiclase.

    \item \textbf{Comparación controlada de arquitecturas:} Evaluación sistemática de
cuatro arquitecturas CNN preentrenadas (MobileNetV2, ResNet50V2, InceptionV3, VGG19) bajo
condiciones experimentales idénticas.

    \item \textbf{Validación externa rigurosa:} Evaluación del mejor modelo en tres
datasets externos independientes (MESSIDOR-2, IDRiD, FOSCAL) con diferentes
características de adquisición y distribución de clases.

    \item \textbf{Combinación de técnicas de regularización:} Integración de bloques
Multi-Head Attention, capas GaussianNoise y técnicas tradicionales (Dropout, L2, early
stopping) para mejorar generalización.

    \item \textbf{Enfoque dual de clasificación:} Desarrollo y evaluación tanto de
clasificadores binarios (optimizados para tamizaje) como multiclase (para gradación
clínica), permitiendo su aplicación en diferentes escenarios clínicos.
\end{itemize}
